% Layout settings
\setcounter{secnumdepth}{3}
\setcounter{tocdepth}{3}

% Global texts
\newcommand{\FIITuniversity}[0] {Slovenská technická univerzita v Bratislave}
\newcommand{\FIITfaculty}[0] {Fakulta informatiky a informačných technológií}
\newcommand{\FIITthesisSlovak}[0] {Bakalárska práca}
\newcommand{\FIITthesis}[0] {Bachelor Thesis}
\newcommand{\FIITtitle}[0] {Názov práce anglicky}
\newcommand{\FIITtitleSlovak}[0] {Názov práce po slovensky}
\newcommand{\FIITauthor}[0] {Meno autora}
\newcommand{\FIITsupervisor}[0] {Meno vedúceho práce}
\newcommand{\FIITevidenceNumber}[0] {FIIT-123456789}
\newcommand{\FIITdate}[0] {Máj 2025}
\newcommand{\FIITassignmentDate}[0] {18.04.2025}
\newcommand{\FIITstudyProgram}[0] {Informatics}
\newcommand{\FIITstudyProgramSlovak}[0] {Informatika}
\newcommand{\FIITdegreeCourse}[0] {2511V00 Applied informatics}
\newcommand{\FIITinstitute}[0] {Institute of Informatics and Software Engineering, FIIT STU, Bratislava}
\newcommand{\FIITsignPlace}[0] {V Bratislave, }
\newcommand{\FIITsignDate}[0] {12.05.2025}
\newcommand{\FIITstudyField}[0]{Computer Science}
\newcommand{\FIITstudyFieldSlovak}[0]{Informatika}
\newcommand{\FIITtrainingWorkplace}[0]{Institute of Computer Engineering and Applied Informatics}
\newcommand{\FIITtrainingWorkplaceSlovak}[0]{Ústav počítačového inžinierstva a aplikovanej informatiky}
\documentclass[12pt,a4paper,oneside,openright,slovak]{article}
\renewcommand{\contentsname}{Obsah}
\usepackage[final]{pdfpages}
\usepackage[
    pdfauthor={""},
    pdftitle={""},
    pdfsubject={""},
    pdfkeywords={""},
    hidelinks
]{hyperref}
\usepackage[utf8]{inputenc}
\usepackage[T1]{fontenc}
\usepackage[a4paper]{geometry}
\usepackage[parfill]{parskip}
\usepackage{xcolor}
\usepackage{enumitem}
\usepackage{calc}
\usepackage{graphicx}
\usepackage{float}
\usepackage{longtable}
\usepackage{setspace}
\usepackage{tabularx}
\usepackage{fancyhdr}
\usepackage[backend=bibtex,sorting=none]{biblatex}
\usepackage{listing}
\usepackage{lscape}
\usepackage{afterpage}
\usepackage{bera}
\usepackage{listings}
\usepackage{lipsum}
\usepackage{lmodern}
\usepackage{fancyvrb}
\usepackage{tocloft}
\usepackage{comment}
\usepackage[most]{tcolorbox}
\usepackage[main=english]{babel}
\usepackage{makecell}
\usepackage{etoolbox}
\usepackage{caption}
\usepackage{titlesec}
\usepackage{multirow}
\usepackage{siunitx}
\usepackage{booktabs}
\usepackage{arydshln}
\usepackage{colortbl}
\usepackage{titlesec}
\usepackage{diagbox}
\usepackage{listings}
\usepackage{geometry}
\usepackage{longtable}
\usepackage{cleveref}
% \usepackage{newclude}
\usepackage{slovak}

\setcounter{secnumdepth}{3}
\setcounter{tocdepth}{3}

\newtcolorbox{todobox}{colback=orange!5!white,colframe=orange!75!black,fonttitle=\bfseries,title=TODO}


\bibliography{references}

% Page design
\pagestyle{fancy}
\lhead{\nouppercase{\leftmark}}
\chead{}
\rhead{}
\lfoot{}
\cfoot{\thepage}
\rfoot{}
 
\begin{document}
\pagenumbering{roman}
\begin{center}
\thispagestyle{empty}

\FIITuniversity
\par\end{center}{\Large \par}

\begin{center}
{\Large \FIITfaculty}
\par\end{center}{\Large \par}

\smallskip{}

\begin{center}
\FIITevidenceNumber
\par\end{center}
\vfill{}

\begin{center}
{\Large \FIITauthor}
\par\end{center}{\Large \par}

\medskip{}

\begin{center}
{\LARGE \FIITtitleSlovak}
\par\end{center}{\huge \par}

\medskip{}

\begin{center}
{\Large \FIITthesisSlovak}
\par\end{center}{\Large \par}

\vfill{}



Vedúci práce: \FIITsupervisor \\

\medskip{}

\FIITdate

%%%%%%%%%%%%%%%%%%%%%%%%%%%%%%%%%%%%%%%%%%%%%%%%%%%%%%%%%%%%%%%%%%%%%%%%%%%%%%%%%%%%%%%%
%%
%% Thesis title page
%%
%%%%%%%%%%%%%%%%%%%%%%%%%%%%%%%%%%%%%%%%%%%%%%%%%%%%%%%%%%%%%%%%%%%%%%%%%%%%%%%%%%%%%%%%
\newpage
\thispagestyle{empty}
\mbox{}
\newpage
\begin{center}
\thispagestyle{empty}

\FIITuniversity
\par\end{center}{\Large \par}

\begin{center}
{\Large \FIITfaculty}
\par\end{center}{\Large \par}

\smallskip{}

\begin{center}
\FIITevidenceNumber
\par\end{center}
\vfill{}

\begin{center}
{\Large \FIITauthor}
\par\end{center}{\Large \par}

\medskip{}

\begin{center}
{\LARGE \FIITtitleSlovak}
\par\end{center}{\huge \par}

\medskip{}

\begin{center}
{\Large \FIITthesisSlovak}
\par\end{center}{\Large \par}

\vfill{}


Študijný program: \FIITstudyProgramSlovak \\
Študijný odbor: \FIITstudyFieldSlovak \\
Miesto vypracovania: \FIITtrainingWorkplaceSlovak \\
Vedúci práce: \FIITsupervisor \\

\medskip{}

\FIITdate

%%%%%%%%%%%%%%%%%%%%%%%%%%%%%%%%%%%%%%%%%%%%%%%%%%%%%%%%%%%%%%%%%%%%%%%%%%%%%%%%%%%%%%%%
%%
%% Thesis title page
%%
%%%%%%%%%%%%%%%%%%%%%%%%%%%%%%%%%%%%%%%%%%%%%%%%%%%%%%%%%%%%%%%%%%%%%%%%%%%%%%%%%%%%%%%%
\newpage
\thispagestyle{empty}
\mbox{}
% \pagebreak
% \begin{center}
\thispagestyle{empty}

\FIITuniversity
\par\end{center}{\Large \par}

\begin{center}
{\Large \FIITfaculty}
\par\end{center}{\Large \par}

\smallskip{}

\begin{center}
\FIITevidenceNumber
\par\end{center}
\vfill{}

\begin{center}
{\Large \FIITauthor}
\par\end{center}{\Large \par}

\medskip{}

\begin{center}
{\LARGE \FIITtitleSlovak}
\par\end{center}{\huge \par}

\medskip{}

\begin{center}
{\Large \FIITthesisSlovak}
\par\end{center}{\Large \par}

\vfill{}


% % Study program: \FIITstudyProgram

% %Degree course: \FIITdegreeCourse

% Place of elaboration: \FIITinstitute

% Study programme: \FIITstudyProgram
% Supervisor: \FIITsupervisor

% \medskip{}

% \FIITdate
% \pagebreak


\thispagestyle{empty}
Vložte zadanie práce


% Bring your own PDF

\newpage
\thispagestyle{empty}
\mbox{}
\newpage
\thispagestyle{empty}

\vspace*{\fill}
{\Large\textbf{Čestné vyhlásenie}}

Čestne vyhlasujem, že som túto prácu vypracoval(a) samostatne, na základe konzultácií \textbf{s XY } a s použitím uvedenej literatúry.


\vspace{2cm}

\begin{tabular}{@{}p{0.5\textwidth}@{}p{0.5\textwidth}@{}}
    \FIITsignPlace \FIITsignDate & 
    \hfill
    \begin{minipage}{5cm}
        \centering
        \hrulefill\\
        \smallskip
        \FIITauthor
    \end{minipage}
\end{tabular}

% \newpage
% \thispagestyle{empty}
% \mbox{}
% \newpage

\thispagestyle{empty}
{\Large\textbf{Poďakovanie}}\\

Poďakovanie ....

%%%%%%%%%%%%%%%%%%%%%%%%%%%%%%%%%%%%%%%%%%%%%%%%%%%%%%%%%%%%%%%%%%%%%%%%%%%%%%%%%%%%%%%%
%%
%% Slovak annotation
%%
%%%%%%%%%%%%%%%%%%%%%%%%%%%%%%%%%%%%%%%%%%%%%%%%%%%%%%%%%%%%%%%%%%%%%%%%%%%%%%%%%%%%%%%%
\newpage
\thispagestyle{empty}
\begin{flushleft}
\begin{Large}
\textbf{Anotácia} \\
\end{Large}
Slovenská technická univerzita v Bratislave \\
FAKULTA INFORMATIKY A INFORMAČNÝCH TECHNOLÓGIÍ\vspace{0.6cm}
\noindent
	\textrm{
	    \begin{normalsize}
	    \begin{tabular}{@{}l>{\raggedright\arraybackslash}p{10cm}@{}}
			Študijný program: & \FIITstudyProgramSlovak\\
			Autor: & \FIITauthor \\
			Bakalárska práca: & \FIITtitleSlovak \\
			Vedúci bakalárskej práce: & \FIITsupervisor \\
			August,\ 2025
	    \end{tabular}
	  	\end{normalsize}
	}
\noindent
\end{flushleft}
Stručná charakteristika zadania bakalárskeho projektu ale predovšetkým výsledkov bakalárskeho projektu v slovenskom jazyku každá v rozsahu max. 1 strany A4 (hlavička + cca 150-200 slov).




\thispagestyle{empty}
\mbox{}
%%%%%%%%%%%%%%%%%%%%%%%%%%%%%%%%%%%%%%%%%%%%%%%%%%%%%%%%%%%%%%%%%%%%%%%%%%%%%%%%%%%%%%%%
%%
%% English annotation
%%
%%%%%%%%%%%%%%%%%%%%%%%%%%%%%%%%%%%%%%%%%%%%%%%%%%%%%%%%%%%%%%%%%%%%%%%%%%%%%%%%%%%%%%%%
\newpage
\thispagestyle{empty}
\begin{flushleft}
\begin{Large}
\textbf{Annotation} \\
\end{Large}
Slovak University of Technology Bratislava \\
FACULTY OF INFORMATICS AND INFORMATION TECHNOLOGIES\vspace{0.6cm}
\noindent
	\textrm{
	    \begin{normalsize}
	    \begin{tabular}{@{}l>{\raggedright\arraybackslash}p{10cm}@{}}
			Study program: & \FIITstudyProgram\\ 
			Author: & \FIITauthor \\
			Bachelor’s thesis: & \FIITtitle \\
			Supervisor: & \FIITsupervisor \\
			August,\ 2025
	    \end{tabular}
	  	\end{normalsize}
	}
\noindent
\end{flushleft}

A brief description of the Bachelor's Project Assignment and, above all, the results of the Bachelor's Project in English, with a maximum of one A4 page each (header + approx. 150-200 words). 



\newpage
\thispagestyle{empty}
\mbox{}
\newpage

\setlength{\cftbeforesecskip}{3pt}

% Table of contents
\tableofcontents


\begin{refsegment}

\thispagestyle{empty}

\pagestyle{fancy}
\titleformat{\section}{\Large\bfseries}{\thesection}{1em}{}
\clearpage
\pagenumbering{arabic}
\clearpage
\section{Technický Abstrakt}

Technický abstrakt je stručný prehľad vašej práce, ktorý ju zasadzuje do širšieho odborného kontextu študovaného odboru. Jeho rozsah býva obmedzený na približne 250 slov a mal by obsahovať jasné a výstižné predstavenie riešenej témy, jej významu v rámci danej domény, ako aj základné ciele a prístup k riešeniu. Rozsah: cca 250 slov. Nesmie presiahnuť 1 stranu.

\clearpage
\section{Laický súhrn}

Nakoľko je technický abstrakt určený profesionálom, laický súhrn by mal obsahovať rovnaký obsah pochopiteľný kýmkoľvek bez profesionálneho vzdelania. Predstavte si, že sa vás na rodinnej oslave niekto spýta čo robíte na bakalársku prácu a vy mu musíte vysvetliť podstatu vašej práce. Toto je zhruba úroveň ktorú očakávame od tejto kapitoly. Rozsah: cca 250 slov. Nesmie presiahnuť 1 stranu


\clearpage
\section{Úvod}
Úvod do problematiky by mal obsahovať základné informácie o projekte a prehľad podobných projektov ktoré riešia daný problém. Samotný úvod by mal byť všeobecnejší, ale nechceme aby obsahoval všeobecné frázy generované veľkými jazykovými modelmi. Úvod by mal predstavovať vaše vlastné chápanie širšej domény a problémoch, ktoré sú v nej riešené na úrovni state-of-the-art. Rozsah do 1 jednej strany. 


\clearpage
\section{Definovanie problému a jeho riešenie}

\subsection{Definovanie problému}
Vaša bakalárska práca rieši jeden konkrétny problém, nie širokú škálu problémov s nejasným pozadím. Očakávame preto, že ste schopní problém pomenovať v krátkej kapitole a konkretizovať riešenia, ktoré sa snažíte v práci implementovať a (po prípade) porovnať. Rozsah: 1-1a ½ strany
\cite{key}
\subsection{Prehľad podobných riešení}
Prezentujte základný literárny rešerš ohľadne podobných riešení vášho problému. Bol tento problém riešený inými ľuďmi alebo v iných oblastiach? V tomto smere očakávame základné zvládnutie technickej literatúry, nie web stránok ani informácií na úrovni blogov. Vyberte a vyhodnoťte odbornú literatúru a zhrňte najvhodnejšie riešenia, porovnajte ich silné stránky a obmedzenia.  Rozsah: 1-1 a ½ strany

\subsection{Prehľad riešenia}
Stručne vysvetlite riešenie problému. Študenti by mali byť schopní odôvodniť svoj výber použitých techník na riešenie daného problému, pričom by mali rozpoznať ich obmedzenia a interdisciplinárny dosah. Rozsah: 2-3 strany

\subsection{Posúdenie rizík}
Pred tým ako sa pustíte do projektu snažte sa vypracovať základnú SWOT analýzu kde pomenujete silné a slabé stránky projektu. Neočakávame v tomto kroku sofistikovanú analýzu podľa industriálnych štandardov, ale základné pochopenie problémov, ktoré váš môžu postretnúť počas vypracovania bakalárskej práce. Pre zistené rizika navrhnite stratégie pre ich riešenie alebo zmiernenie ich dopadu. Rozsah do 1 strany



\subsection{Experimentálna reprodukovateľnosť a integrácia}
Nech už vykonávate akúkoľvek činnosť, dbajte na to, aby vaše kroky boli transparentné a časovo reprodukovateľné. Táto kapitola by mala jednoznačne identifikovať proaktívne opatrenia, ktoré ste prijali s cieľom zabezpečiť, že výsledná práca bude v súlade s princípmi otvorenej vedy a reprodukovateľného výskumu. Vzhľadom na rôznorodosť tejto oblasti, zahŕňajúcej široké spektrum prístupov — od kontajnerizácie, dôslednej dokumentácie kódu, až po sprístupnenie všetkých komponentov projektu vo verejných repozitároch — je dôležité uviesť aj obmedzenia vyplývajúce z dohôd o mlčanlivosti (NDA), ktoré môžu brániť úplnému zverejneniu niektorých častí. Rozsah: 3 strany



\subsection{Udržateľnosť a vplyv na životné prostredie}
Nech už pracujete na akomkoľvek projekte, pristupujte k nemu tak, aby mal dlhodobú hodnotu. Mnohé študentské práce sú navrhované s cieľom splniť požiadavky bakalárskeho štúdia, no po niekoľkých rokoch bývajú nefunkčné, neudržiavané a prakticky nepoužiteľné. V tejto kapitole preto navrhnite stratégiu, ako ste zabezpečili dlhodobú udržateľnosť vášho riešenia. Popíšte, aké kroky ste podnikli na minimalizáciu potreby zásahov v budúcnosti, ako plánujete dokumentovať a archivovať projekt, alebo akým spôsobom budú vaše ML modely pripravené na rozšírenie bez nutnosti opätovného trénovania pri dostupnosti nových dát. Špeciálne v poslednom prípade zvážte environmentálny dopad ďalších zásahov pri udržiavaní vášho riešenia. Rozsah: 1-2 strany
\subsection{Zamestnateľnosť}
Predpokladáme, že výber projektu vychádzal z vašich osobných záujmov a motivácie rozvíjať sa v oblasti, ktorá vás skutočne zaujíma. Počas prvého semestra by ste už mali vedieť zhodnotiť, ako práca na projekte prispieva k vášmu odbornému rastu a akým spôsobom zvyšuje vašu konkurencieschopnosť na trhu práce. Táto kapitola by preto mala reflektovať potenciál projektu z hľadiska kariérneho rozvoja, nadobudnutých zručností a jeho prepojenia s aktuálnymi požiadavkami praxe. Rozsah: ½ strany
\subsection{Tímová práca, diverzita a inklúzia}
Vaša rola v projekte bude vždy definovaná ako IT expert, no rovnako dôležité je, aby ste dokázali vnímať a rešpektovať perspektívu kolegov bez technického zázemia. V tejto kapitole popíšte, ako ste komunikovali, vysvetľovali technické riešenia a navrhovali postupy tak, aby boli zrozumiteľné a prakticky využiteľné aj pre neinformatikov, s ktorými ste spolupracovali. Zároveň uveďte, aké procesy ste použili pri rozdelení rolí v tíme, aký systém projektového riadenia ste zvolili, ako prebiehali spoločné porady a akým spôsobom ste dospeli ku konsenzu v situáciách, kde sa názory jednotlivých členov tímu rozchádzali. Rozsah: ½ strany

\clearpage
\section{Zhodnotenie}
Záverečná kapitola by mala byť vecná, koncentrovaná a bez opakovania už uvedených skutočností. Má zhrnúť kľúčové aspekty práce zohľadňujúc vplyv na spoločnosť a inžiniersku prax.  Reflektujte projekt z nadhľadu: aký mal prínos pre váš odborný rast, ako obohatil vaše pracovné prostredie a akú hodnotu prináša jeho výsledok cieľovému používateľovi. Zároveň sa zamyslite nad širším dopadom vašej práce — či už ide o prínos k spoločenským výzvam, podporu interdisciplinárnej spolupráce, alebo posun v inžinierskej praxi v danom odbore. Cieľom je ukázať, že projekt má presah nad rámec akademického zadania a dokáže generovať hodnotu aj v reálnom svete. Rozsah: 2 strany
% After your Conclusions chapter but before the end of the document
\clearpage
% Bibliography
\addcontentsline{toc}{section}{Referencie}
Obsahuje zoznam použitej literatúry (maximálne 70) musí obsahovať všetky náležitosti, tak ako uvádza norma STN ISO 690. V texte bakalárskeho projektu sa musia nachádzať odvolávky na každú citovanú literatúru. 

\printbibliography[title=Referencie,segment=\therefsegment]

% For chapters/14_Scientific_part.tex:
\section{Vedecká časť práce}
Použite formát typický pre vedecké články:  3500 slov, maximálne 7 obrázkov, jeden z dvoch bežných formátov (abstrakt, laický abstrakt, úvod, výsledky, diskusia, metódy, doplnkový materiál a údaje) alebo (abstrakt, laický abstrakt, úvod, metódy, výsledky, diskusia, doplnkový materiál a údaje). Štruktúra a rozsah sú iba odporúčané a mali by sa určiť po konzultácii s vedúcim práce v závislosti od riešenej problematiky.

Od študentov sa očakáva, že preukážu svoju schopnosť efektívne komunikovať riešenia zložitých problémov technickému publiku. V tejto časti nie je povolené používať obsah generovaný umelou inteligenciou.

Formátovanie vedeckého článku je podľa aktuálnej šablóny IIT.SRC, ktorá je dostupná na stránke https://iitsrc.sk/submission/.


% Include PDF with no decorations



% List of Figures
% \addcontentsline{toc}{section}{List of Figures}
% \listoffigures

\thispagestyle{empty}
\mbox{}

\appendix
\captionsetup[figure]{list=no}
\captionsetup[table]{list=no}

\begin{titlepage}
    \centering
    \vspace*{\fill}
    {\Huge\bfseries Prílohy\par}
    \vspace*{\fill}
\end{titlepage}

% Start Appendix A
\clearpage
\renewcommand{\thepage}{A-\arabic{page}}
\setcounter{page}{1}
\section{Prílohy}
Do tejto časti patria obrazové prílohy, technická dokumentácia kódu, dockerizované aplikácie a ďalšie relevantné výstupy, ktoré dopĺňajú projekt a podporujú jeho transparentnosť, reprodukovateľnosť a praktickú využiteľnosť. Rozsah materiálu musí byť dostačujúci pre úspešné reprodukovanie popísaných experimentov.
Súčasťou je aj plán práce riešenia projektu spolu so študentovým vyjadrením k jeho plneniu. V rámci práce je povinné uviesť plán práce.
V rámci tejto časti je tiež popísaná štruktúra a účel odovzdaných súborov záverečnej práce v prehľadnej forme ako sú napr. zdrojové kódy, výsledky, grafy atď.



% Start Appendix B
\clearpage
\renewcommand{\thepage}{B-\arabic{page}}
\setcounter{page}{1}
\section {Plán práce}

\subsection*{Počiatočný výskum pred začiatkom semestra}
\begin{table}[H]
\centering
\begin{tabularx}{\textwidth}{lX}
\textbf{Mesiac} & \textbf{Aktivita} \\
\hline
Jún & Interdum et malesuada fames ac ante ipsum primis in faucibus. Pellentesque sed auctor est. \\
 & Interdum et malesuada fames ac ante ipsum primis in faucibus. Pellentesque sed auctor est. \\
 
\hline
Júl & In bibendum quis ligula ac viverra. \\
 & In bibendum quis ligula ac viverra. \\
 & In bibendum quis ligula ac viverra. \\
\hline
\end{tabularx}
\caption{Plán práce pred začiatkom zimného semestra}
\label{tab:summer-plan}
\end{table}

\paragraph{Vyhodnotenie plánu prace pred začatím zimného semestra}
Proin et fermentum diam, et sollicitudin enim. Duis erat nunc, rhoncus sit amet eleifend ut, semper pharetra ligula. In facilisis augue nulla, quis blandit mauris accumsan et. Interdum et malesuada fames ac ante ipsum primis in faucibus. Pellentesque sed auctor est. Cras sed est eu tellus rhoncus rhoncus non eget nisl. In bibendum quis ligula ac viverra. Integer sapien elit, dapibus eu justo nec, fringilla iaculis leo. Sed condimentum sollicitudin leo, nec facilisis lectus lobortis sed. Mauris euismod congue neque, sed sodales velit pellentesque vitae. Ut rutrum eget velit ac convallis. Cras eu tincidunt nulla. Nulla interdum rhoncus ultricies. Aliquam tortor elit, fermentum vel eleifend cursus, porttitor nec ante.

\subsection*{Bakalárska práca I. - zimný semester}
\begin{table}[H]
\centering
\begin{tabularx}{\textwidth}{lX}
\textbf{Týždeň} & \textbf{Aktivita} \\
\hline
1. – 2. týždeň & Proin et fermentum diam, et sollicitudin enim. \\
3. – 4. týždeň & Pharetra ligula. In facilisis augue nulla. \\
\hline
\end{tabularx}
\caption{Plán práce pre zimný semester}
\label{tab:winter-plan}
\end{table}

\paragraph{Vyhodnotenie plánu prácu zimného semestra}
Proin et fermentum diam, et sollicitudin enim. Duis erat nunc, rhoncus sit amet eleifend ut, semper pharetra ligula. In facilisis augue nulla, quis blandit mauris accumsan et. Interdum et malesuada fames ac ante ipsum primis in faucibus. Pellentesque sed auctor est. Cras sed est eu tellus rhoncus rhoncus non eget nisl. In bibendum quis ligula ac viverra. Integer sapien elit, dapibus eu justo nec, fringilla iaculis leo. Sed condimentum sollicitudin leo, nec facilisis lectus lobortis sed. Mauris euismod congue neque, sed sodales velit pellentesque vitae. Ut rutrum eget velit ac convallis. Cras eu tincidunt nulla. Nulla interdum rhoncus ultricies. Aliquam tortor elit, fermentum vel eleifend cursus, porttitor nec ante.

\subsection*{Bakalárska práca II. - letný semester}
\begin{table}[H]
\centering
\begin{tabularx}{\textwidth}{lX}
\textbf{Week} & \textbf{Activity} \\
\hline
1. týždeň & Lorem ipsum dolor sit amet, \\
2. týždeň & Lorem ipsum dolor sit amet, \\
 \\

\hline
\end{tabularx}
\caption{Plán práce pre letný semester}
\label{tab:summer2-plan}
\end{table}

\paragraph{Vyhodnotenie plánu prácu letného semestra}
Proin et fermentum diam, et sollicitudin enim. Duis erat nunc, rhoncus sit amet eleifend ut, semper pharetra ligula. In facilisis augue nulla, quis blandit mauris accumsan et. Interdum et malesuada fames ac ante ipsum primis in faucibus. Pellentesque sed auctor est. Cras sed est eu tellus rhoncus rhoncus non eget nisl. In bibendum quis ligula ac viverra. Integer sapien elit, dapibus eu justo nec, fringilla iaculis leo. Sed condimentum sollicitudin leo, nec facilisis lectus lobortis sed. Mauris euismod congue neque, sed sodales velit pellentesque vitae. Ut rutrum eget velit ac convallis. Cras eu tincidunt nulla. Nulla interdum rhoncus ultricies. Aliquam tortor elit, fermentum vel eleifend cursus, porttitor nec ante.

% Start Appendix C
\clearpage
\renewcommand{\thepage}{C-\arabic{page}}
\setcounter{page}{1}
\input{attachments/appendix_C_technical_documentation}





% Start Appendix D
\clearpage
\renewcommand{\thepage}{D-\arabic{page}}
\setcounter{page}{1}
\clearpage
\section{Obsah digitálneho média}
% Force page numbering format with multiple approaches
\renewcommand{\thepage}{D-\arabic{page}}
\setcounter{page}{1}
% Additional forceful page style settings
\makeatletter
\def\ps@specialappendix{%
  \def\@oddhead{}%
  \def\@evenhead{}%
  \def\@oddfoot{\hfil D-\arabic{page}\hfil}%
  \def\@evenfoot{\hfil D-\arabic{page}\hfil}%
}
\makeatother
\pagestyle{specialappendix}
\thispagestyle{specialappendix}
Thesis registration number in the information system: \FIITevidenceNumber.

The attached digital medium contains the complete source code, configuration files, and scripts for this thesis. The ZIP archive is organized as follows:

\begin{tabular}{@{}p{0.45\textwidth}p{0.55\textwidth}@{}}
% No header, completely removed the header row
\ttfamily README.md & Project overview and setup instructions \\[0.5em]
\ttfamily applications/ & Source code for all applications \\
\ttfamily \hspace{0.5cm} operators/ & Carbon-aware applications \\

\end{tabular}
% Apply the special page style to next page too
\clearpage
\thispagestyle{specialappendix}

\end{refsegment}
\end{document}